\documentclass[a4paper,11pt]{article}
\usepackage[utf8]{inputenc}
\usepackage{fancyvrb}
\usepackage[margin=2.5cm]{geometry}
\usepackage[spanish]{babel}
\usepackage{graphicx}
\usepackage{tabularx}
\usepackage{multirow}
\usepackage{booktabs}
\def\tablename{Tabla}

\title{Manual del código de MATLAB/Octave de Calculo Matricial de Estructuras de NUdos R\'igidos (CMENUR)}
\author{Jacob David Rodr\'iguez Bord\'on}
\date{Marzo 2022}

\begin{document}

\maketitle

Este código de MATLAB/Octave permite analizar en régimen estático estructuras barras de nudos rígidos bidimensionales, esto es, cálculo de desplazamientos, giros, esfuerzos axiles, esfuerzos cortantes y momentos flectores. La implementación del mismo tiene un objetivo educativo, y no profesional.

Para la definición del modelo se debe escribir un fichero de entrada de texto plano con el siguiente formato:
\begin{Verbatim}[frame=single,fontsize=\small]
<número de materiales>
<identificador> <módulo de elasticidad longitudinal>
...

<número de secciones>
<identificador> <área> <inercia>
...

<número de nodos>
<identificador> <x> <y>
...

<número de elementos>
<identificador> <id. nodo i> <id. nodo j> <id. sección> <id. material>
...

<número de apoyos>
<identificador> <id. nodo> <tipo de apoyo> [valores adicionales]
...

<número de cargas nodales>
<identificador> <id. nodo> <Fx> <Fy> <M>
...

<número de cargas transversales>
<identificador> <id. elemento> <q> 
...

\end{Verbatim}
donde la selección del tipo de apoyo y los \texttt{[valores adicionales]} están definidos en la Tabla \ref{tab:valores}. Tómense los casos \texttt{ej1.dat} y \texttt{ej2.dat} en la carpeta \texttt{examples/} como ejemplos.

Para realizar el análisis, basta con situar la carpeta de trabajo en la raíz del programa, y luego ejecutar la función \texttt{cmenur} con un único argumento de entrada indicando la ruta hacia el fichero de entrada, por ejemplo (ruta al archivo para sistemas Linux):
\begin{Verbatim}[frame=single,fontsize=\small]
> modelo_analizado = cmenur('examples/ej1.dat');
\end{Verbatim}
La función devuelve una variable (una estructura de datos), llamada en este caso \texttt{modelo\_analizado}, que contiene no sólo los datos del modelo sino los resultados del análisis. Una descripción de esta variable puede verse en \texttt{leer\_datos.m}. Además, se generan tres archivos de salida, en este caso de ejemplo: \texttt{ej1.dat.u.txt}, \texttt{ej1.dat.F.txt} y \texttt{ej1.dat.R.txt}; en donde es escriben los resultados de desplazamientos y giros, esfuerzos en las barras, y reacciones sobre los apoyos (no las reacciones de los apoyos sobre la estructura).

\begin{table}
\begin{center}
\renewcommand{\arraystretch}{1.5}
\begin{tabular}{c p{4cm} c l}
\toprule
\textbf{Tipo de} & \multirow{2}{*}{\textbf{Descripci\'on}} & \textbf{Representaci\'on} & \textbf{Valores}     \\
\textbf{apoyo}   &                                         & \textbf{gr\'afica}        & \textbf{adicionales} \\
\midrule
1 & Apoyo r\'igido fijo & 
\multicolumn{1}{m{2cm}}{\centering\includegraphics{figures/apoyo_fijo.eps}} & 
\texttt{<b>} \\
\hline
2 & Apoyo r\'igido carro horizontal & 
\multicolumn{1}{m{2cm}}{\centering\includegraphics{figures/apoyo_carro_horizontal.eps}} &
\texttt{<b>} \\
\hline
3 & Apoyo r\'igido carro vertical & 
\multicolumn{1}{m{2cm}}{\centering\includegraphics{figures/apoyo_carro_vertical.eps}} & 
\texttt{<b>}\\
\hline
4 & Apoyo r\'igido carro inclinado & 
\multicolumn{1}{m{2cm}}{\centering\includegraphics{figures/apoyo_inclinado.eps}} & 
\texttt{<a>} \\
\hline
5 & Apoyo tipo fijo empotrado & 
\multicolumn{1}{m{2cm}}{\centering\includegraphics{figures/apoyo_empotramiento.eps}} & 
\texttt{<b>} \\
\hline
6 & Apoyo tipo carro horizontal empotrado & 
\multicolumn{1}{m{2cm}}{\centering\includegraphics{figures/apoyo_empotrado_carro_horizontal.eps}} & 
\texttt{<b>} \\
\hline
7 & Apoyo tipo carro vertical empotrado & 
\multicolumn{1}{m{2cm}}{\centering\includegraphics{figures/apoyo_empotrado_carro_vertical.eps}} & 
\texttt{<b>} \\
\hline
8 & Apoyo tipo carro inclinado empotrado & 
\multicolumn{1}{m{2cm}}{\centering\includegraphics{figures/apoyo_empotrado_carro_inclinado.eps}} & 
\texttt{<a>} \\
\hline
9 & Apoyo el\'astico & 
\multicolumn{1}{m{2cm}}{\centering\includegraphics{figures/apoyo_elastico.eps}} & 
\texttt{<Kx>}  \texttt{<Ky>} \texttt{<Kr>} \texttt{<bx>} \texttt{<by>} \texttt{<br>} \\
\bottomrule
\end{tabular}
\end{center}
\caption{Valores adicionales para definir los apoyos. Nota: \texttt{<a>} es el \'angulo $\alpha$ de giro de los apoyos inclinados, \texttt{<b>} es el \'angulo $\beta$ usado para girar la representacion gr\'afica mostrada, los \'angulos han de introducirse en grados)}
\label{tab:valores}
\end{table}

\end{document}

